% !PRINTED 116
\weberpage{116}
Setzt man ferner $u = 0$ in (7), (8), (9), \S\,33, benutzt alsdann die Relationen (10) und (20) und zieht die Quadratwurzel, so findet sich

\weq{21}{\begin{aligned}
f(n\,\omega)\,\eta(\omega)^{\frac{n-1}{2}} &= f(\omega) \prod_{1,\,\frac{n-1}{2}}^{\nu} \vartheta_{00}\!\left(\frac{\nu}{n}\right), \\
f_1(n\,\omega)\,\eta(\omega)^{\frac{n-1}{2}} &= f_1(\omega) \prod_{1,\,\frac{n-1}{2}}^{\nu} \vartheta_{01}\!\left(\frac{\nu}{n}\right), \\
f_2(n\,\omega)\,\eta(\omega)^{\frac{n-1}{2}} &= f_2(\omega) \prod_{1,\,\frac{n-1}{2}}^{\nu} \vartheta_{10}\!\left(\frac{\nu}{n}\right).
\end{aligned}}

Die Vorzeichen ergeben sich hier aus der Annahme eines unendlich kleinen $q$.

\begin{center}
\textbf{\S\,35. \gesperrt{Die Weierstrasssche $\sigma$-Funktion.}}
\end{center}

Durch die allgemeine lineare Substitution
\[
\begin{pmatrix} \alpha, & \beta \\ \gamma, & \delta \end{pmatrix}, \qquad \alpha\delta - \beta\gamma = 1,
\]
angewandt auf die Variablen $\omega_1$, $\omega_2$, geht jede $t$-Funktion wieder in eine $t$-Funktion über, wobei die Charakteristik sich geändert hat. Ist $(g_1, g_2)$ die Charakteristik der ursprünglichen, $(g_1', g_2')$ die der transformierten Funktion, so ist nach \S\,27, (11)
\[
(g_1, g_2) = (\delta g_1' - \beta g_2' - \beta\delta,\ -\gamma g_1' + \alpha g_2' - \alpha\gamma).
\]

Löst man die beiden darin enthaltenen linearen Gleichungen für $g_1'$, $g_2'$ auf und beachtet, daß $\alpha\beta(\gamma+\delta+1)$, $\gamma\delta(\alpha+\beta+1)$ notwendig gerade Zahlen sind, und daß eine Charakteristik sich nicht ändert, wenn sich ihre Elemente um Vielfache von 2 ändern, so kann man dafür auch setzen:

\weq{1}{(g_1', g_2') = (\alpha g_1 + \beta g_2 + \alpha\beta,\ \gamma g_1 + \delta g_2 + \gamma\delta).}

Die $\vartheta$-Funktionen gehen also, abgesehen von Exponentialfaktoren, ineinander über. Von Wichtigkeit ist es aber, eine Funktion zu bilden, die den linearen Transformationen gegenüber \gesperrt{absolut invariant} ist, und eine solche Funktion ist die von \gesperrt{Weierstrass} in die Theorie eingeführte $\sigma$-Funktion, zu deren Definition wir jetzt übergehen.

Wenn wir die Formel (1) auf die vier Hauptcharakteristiken. $(0, 0)$, $(0, 1)$, $(1, 0)$, $(1, 1)$ anwenden, so gehen diese der Reihe
% [sic] period after "Hauptcharakteristiken" as printed, where a colon would be expected
%CONT>>
